% =============================================================================
%  CHAPTER 3 — METHODOLOGY (~15-35% of thesis)
% =============================================================================
\section{Methodology}
\label{sec:methodology}

This section describes the design and implementation of the three system
variants developed in this thesis: the RAG system
(Section~\ref{sec:rag-system}), the prompt-only baseline
(Section~\ref{sec:prompt-only}), and the hybrid engine
(Section~\ref{sec:hybrid}). We first present an overview of the overall
architecture and end-to-end pipeline, then detail each component in the
subsequent subsections.

% ─────────────────────────────────────────────────────────────────────────────
\subsection{System Overview}
\label{subsec:system-overview}

The system follows the canonical Retrieval-Augmented Generation (RAG)
architecture introduced by Lewis~et~al.~\cite{lewis2020rag}, in which a
retrieval component supplies relevant external knowledge to a generative
language model at inference time, thereby grounding its outputs in
verifiable sources rather than relying solely on parametric
knowledge~\cite{gao2024rag_survey}. In our implementation, this architecture
is instantiated as a five-stage pipeline specifically tailored to the
Tunisian Baccalaureate mathematics domain. Figure~\ref{fig:system-architecture}
provides a schematic overview of the complete pipeline.

% ── Pipeline diagram (TikZ) ─────────────────────────────────────────────────
\begin{figure}[H]
\centering
\resizebox{\textwidth}{!}{%
\begin{tikzpicture}[
    node distance=1.2cm and 1.8cm,
    stage/.style={rectangle, draw=black!70, fill=blue!8, rounded corners=4pt,
                  minimum width=3.2cm, minimum height=1.1cm, align=center,
                  font=\small\bfseries},
    comp/.style={rectangle, draw=black!50, fill=white, rounded corners=2pt,
                 minimum width=2.6cm, minimum height=0.8cm, align=center,
                 font=\small},
    arrow/.style={->, >=stealth, thick, draw=black!60},
    label/.style={font=\footnotesize\itshape, text=black!60},
]

% ── Stage 1: Data Acquisition ──
\node[stage] (acq) {Stage 1\\Data Acquisition};
\node[comp, below=0.5cm of acq] (scan) {Scanned exams,\\series, cours};

% ── Stage 2: Digitization ──
\node[stage, right=2cm of acq] (dig) {Stage 2\\Digitization};
\node[comp, below=0.5cm of dig] (gemini_ocr) {Gemini multimodal\\OCR $\to$ \LaTeX{}};

% ── Stage 3: Indexing ──
\node[stage, right=2cm of dig] (idx) {Stage 3\\Indexing};
\node[comp, below=0.5cm of idx] (chunk) {Adaptive chunking\\+ BGE-M3 embedding};
\node[comp, below=0.4cm of chunk] (chroma) {ChromaDB\\vector store};

% ── Stage 4: Retrieval ──
\node[stage, right=2cm of idx] (ret) {Stage 4\\Retrieval};
\node[comp, below=0.5cm of ret] (twostage) {Two-stage retrieval\\+ companion fetch};

% ── Stage 5: Generation ──
\node[stage, right=2cm of ret] (gen) {Stage 5\\Generation};
\node[comp, below=0.5cm of gen] (llm) {Gemini LLM\\``Monsieur Tounsi''};

% ── Arrows ──
\draw[arrow] (acq) -- node[above, label] {images/PDF} (dig);
\draw[arrow] (dig) -- node[above, label] {.tex files} (idx);
\draw[arrow] (idx) -- node[above, label] {vectors} (ret);
\draw[arrow] (ret) -- node[above, label] {context} (gen);

% ── Vertical arrows ──
\draw[arrow] (scan) -- (acq);
\draw[arrow] (gemini_ocr) -- (dig);
\draw[arrow] (chunk) -- (idx);
\draw[arrow] (chroma) -- (chunk);
\draw[arrow] (twostage) -- (ret);
\draw[arrow] (llm) -- (gen);

% ── Student query ──
\node[comp, above=0.8cm of ret, fill=green!10, draw=green!50!black]
    (query) {Student query};
\draw[arrow, draw=green!50!black] (query) -- (ret);

% ── Answer output ──
\node[comp, above=0.8cm of gen, fill=red!8, draw=red!50!black]
    (answer) {Grounded answer\\with sources};
\draw[arrow, draw=red!50!black] (gen) -- (answer);

% ── Offline / Online labels ──
\draw[dashed, gray] ([xshift=0.9cm]idx.east |- acq.north) ++(0,1.2)
    -- ++(0,-5.5) node[below, font=\footnotesize\bfseries, text=gray]
    {offline $\mid$ online};

\end{tikzpicture}%
}
\caption{End-to-end pipeline of the RAG system. Stages~1--3 (left of the
dashed line) are executed offline during corpus preparation. Stages~4--5
(right) are executed online at query time. The student's question triggers
a two-stage retrieval over the pre-built vector store, and the retrieved
context is compiled into a prompt for the Gemini language model, which
produces a grounded answer in the Tunisian Baccalaureate redaction style.}
\label{fig:system-architecture}
\end{figure}

\subsubsection{Pipeline Stages}

The five stages of the pipeline are summarized below, with references to the
detailed subsections that follow.

\paragraph{Stage~1: Data Acquisition (Section~\ref{subsec:data-acquisition}).}
The corpus is assembled from three categories of Tunisian Baccalaureate
mathematics material: corrected past official examinations
(\textit{bac\_officiel}), chapter-based exercise series with worked
corrections (\textit{série}), and course notes with theorem statements
(\textit{cours}). A deliberate design decision was to acquire each exercise
\textit{together with its correction} as a single logical unit, preserving
the paired structure that is fundamental to mathematical
pedagogy~\cite{zhao2024financial_chunking, chen2024mog}. The final corpus
comprises 624 source files spanning 20 chapters of the Baccalaureate
mathematics program (Section Mathématiques).

\paragraph{Stage~2: Digitization (Section~\ref{subsec:digitization}).}
Because the source material consists of scanned images and PDF files, a
digitization step converts each document into machine-readable \LaTeX{} text.
We use Google Gemini as a multimodal OCR engine, prompting it to faithfully
transcribe both the natural-language text and the mathematical expressions
while preserving the official Tunisian redaction
style~\cite{wang2023ocrbench, mathpix2024}. The digitization pipeline is
incremental and resume-friendly: previously processed files are detected via
companion \texttt{.tex} files in Google Cloud Storage and skipped on
subsequent runs.

\paragraph{Stage~3: Indexing (Section~\ref{subsec:database-construction}).}
The digitized \LaTeX{} files are normalized, split into chunks using an
\textit{adaptive chunking strategy} that adjusts chunk size based on document
type (1{,}500 characters for corrections, 3{,}000 for course material), and
embedded into dense vectors using the BGE-M3 multilingual embedding
model~\cite{chen2024bge_m3}. Each chunk is annotated with rich metadata
(document type, chapter, year, session, exercise identifier, solution flag)
extracted from the file path structure. The resulting 1{,}422 chunks and
their embeddings are stored in ChromaDB~\cite{chromadb2024}, a persistent
vector database with native metadata filtering support. An incremental
manifest mechanism ensures that only modified files are re-embedded on
subsequent indexing runs.

\paragraph{Stage~4: Two-Stage Retrieval
(Section~\ref{subsec:retrieval-engine}).}
When a student submits a question, the retrieval engine executes a
\textit{two-stage} similarity search over the vector store. The first pass
targets correction-type documents (\texttt{is\_solution=true}), seeking
worked examples whose redaction style the language model can
mimic~\cite{karpukhin2020dpr}. If no sufficiently similar correction is
found (based on L2 distance thresholds calibrated empirically), a second pass
retrieves course material to provide theorem-level backing. A companion
fetch mechanism then locates the exercise statement corresponding to each
retrieved correction, ensuring that the language model receives the complete
exercise--correction pair as context~\cite{gao2024rag_survey}. This
two-stage design reflects the pedagogical distinction between ``knowing the
theory'' and ``knowing how to write a Baccalaureate-style solution,'' and
is a key architectural contribution of this work.

\paragraph{Stage~5: Generation
(Section~\ref{subsec:retrieval-engine}).}
The retrieved documents are compiled into a structured prompt alongside the
student's question and a system persona (``Monsieur Tounsi,'' a veteran
Tunisian mathematics teacher). The prompt enforces anti-hallucination rules,
syllabus boundary constraints, and the formal redaction conventions of the
Tunisian Baccalaureate~\cite{kasneci2023chatgpt_education}. Generation is
performed by Google Gemini via the Vertex AI API, with a low temperature
($T = 0.15$) to favor deterministic mathematical
reasoning~\cite{team2023gemini}. The final answer is accompanied by a
confidence level (high, medium, or low) computed from retrieval quality
metrics, and full source transparency is provided through an expandable
panel listing the retrieved documents and their metadata.

\subsubsection{Shared Infrastructure Across System Variants}

All three system variants evaluated in this thesis—the RAG system, the
prompt-only baseline, and the hybrid engine—share the same underlying
infrastructure:

\begin{itemize}
    \item The same generative model (Google Gemini, accessed through
    Vertex AI) with identical temperature ($T = 0.15$) and maximum output
    token settings (4{,}096 tokens).
    \item The same system persona (``Monsieur Tounsi'') with the same
    syllabus guard, redaction conventions, and anti-hallucination
    constraints.
    \item The same Streamlit-based user interface, providing a bilingual
    chat interface (French and Tunisian Arabic) with source transparency
    and optional debug observability.
\end{itemize}

The \textit{only} difference between the three variants is how each
constructs the context portion of the prompt: the RAG system includes
retrieved documents from the vector store, the prompt-only baseline encodes
the curriculum directly into the system prompt through ten specialized prompt
components, and the hybrid engine dynamically selects between the two
strategies based on retrieval confidence (Section~\ref{sec:hybrid}). This
controlled design ensures that any performance difference observed in the
evaluation (Section~\ref{sec:experiments}) can be attributed to the retrieval
mechanism rather than to confounding factors in the generation pipeline.

% ─────────────────────────────────────────────────────────────────────────────
% Include the existing detailed RAG methods section
% The detailed RAG methods section is in the parent thesis directory.
% When compiling from thesis/, use: pdflatex main.tex
% =============================================================================
%  RAG SYSTEM — detailed methods (included from methodology.tex)
% =============================================================================

\subsection{RAG System}
\label{sec:rag-system}

This section presents the RAG system, which is the central technical
contribution of this thesis. The pipeline has six components: data
acquisition, digitization, database construction, embedding model selection,
two-stage retrieval, and answer generation.


% ─────────────────────────────────────────────────────────────────────────────
\subsubsection{Data Acquisition}
\label{subsec:data-acquisition}

The foundation of any RAG system is the quality of its knowledge
base~\cite{lewis2020rag, gao2024rag_survey}. Research on data quality in
AI systems has shown that downstream performance is often dominated by data
quality rather than model architecture, and that compounding quality issues
from early data stages are the most common cause of
failure~\cite{sambasivan2021everyone}. This guided our approach: we
prioritized corpus quality and structural consistency over size.

No publicly available, machine-readable corpus of Tunisian Bac mathematics
materials existed before this work. The primary sources are printed
examination booklets, handwritten corrections, and scanned course pages
from preparatory-class teachers.

\paragraph{Physical data collection.}
The raw material was collected during a trip to Tunisia, where we obtained
books and course materials from experienced Bac mathematics teachers. For
printed material, a standard scanner or phone camera produced sufficient
quality. For handwritten corrections with dense notation or faded ink, the
corrections were manually rewritten on a tablet (iPad with Apple Pencil)
before scanning---a labor-intensive but necessary investment in data
quality~\cite{sambasivan2021everyone, wang2023ocrbench,
mahdavi2019icdar_math}.

Each exercise was scanned \textit{together with its correction} as a single
image or PDF, rather than separating statements from solutions. Course
material was scanned as continuous documents per chapter.

\paragraph{Corpus structure.}
The material was uploaded to Google Cloud Storage~\cite{gcloud_storage2024}
in a hierarchical structure:

\begin{verbatim}
BacMath_Raw_Data/
  01_Continuite_et_limites/
    Cours/                      -> course material
    series_et_corrections/      -> exercises + corrections
  ...
  20_Coniques/
  Bac_avec_corrections_Images/  -> past exam corrections
\end{verbatim}

This naming convention is machine-readable: the indexing pipeline extracts
document type, chapter, year, session, and exercise number directly from
folder paths using regular expressions. The corpus comprises 624 source
files spanning 20 chapters: 312 corrected exam pages, 250 series exercise
files, and 62 course material files, yielding 1{,}422 chunks after indexing.

\paragraph{Paired exercise--correction acquisition.}
Scanning each exercise together with its correction was a deliberate design
decision. In mathematical pedagogy, the correction is a complete reasoning
chain whose structure depends on the exercise statement. Separating them
would sever this dependency. The retrieval system also benefits: the most
useful retrieval target is a complete worked example, not an orphaned
solution. This aligns with research showing that structure-aware document
processing improves retrieval
quality~\cite{zhao2024financial_chunking, chen2024mog}.


% ─────────────────────────────────────────────────────────────────────────────
\subsubsection{Mathematical Document Digitization}
\label{subsec:digitization}

The next step converts scanned images into machine-readable \LaTeX{}.
Mathematical documents contain two-dimensional structures (fractions,
subscripts, integrals) that standard OCR cannot
handle~\cite{wang2023ocrbench, mahdavi2019icdar_math}.

\paragraph{Tool selection.}
We initially considered Mathpix~\cite{mathpix2024}, a commercial math-OCR
service, but its per-page pricing made it impractical for a thesis prototype.
Instead, we used Gemini via Vertex AI. Benchmarking by Wang et
al.~\cite{wang2023ocrbench} found that Gemini achieved the highest overall
OCR score among multimodal models tested. For our material---primarily
printed or cleanly rewritten---Gemini provided sufficient fidelity at
negligible cost. The choice was driven by budget constraints, not a claim of
superior accuracy over Mathpix.

\paragraph{Pipeline (\texttt{digitize.py}).}
The digitization pipeline: (1)~enumerates all images and PDFs in GCS,
(2)~checks for existing companion \texttt{.tex} files (making the pipeline
resume-friendly), (3)~sends each pending file to Gemini with a structured
prompt requesting faithful \LaTeX{} transcription in the official redaction
style, and (4)~uploads the result back to GCS. The pipeline includes
exponential-backoff retry logic for transient API errors.

\paragraph{Quality.}
Digitization errors propagate through the entire system: a misrecognized
formula gets embedded, retrieved, and injected into the LLM's prompt. The
ground truth evaluation (Section~\ref{subsec:ocr-evaluation}) confirmed an
average CER of 5.4\% and \LaTeX{} math accuracy of 97.1\%---sufficient for
our corpus but not error-free.


% ─────────────────────────────────────────────────────────────────────────────
\subsubsection{Database Construction and Indexing}
\label{subsec:database-construction}

\paragraph{Migration to ChromaDB.}
An early prototype used a JSON-based vector store. This was functional for
initial experiments but could not scale (full memory load per query, no
incremental updates, no metadata filtering). We migrated to
ChromaDB~\cite{chromadb2024}, which provides persistent local storage,
native metadata filtering, and deterministic upsert operations.

\paragraph{Text normalization.}
Before chunking, each \texttt{.tex} file is normalized: CRLF $\to$ LF,
collapsed whitespace, and excess newlines reduced to double newlines.

\paragraph{Adaptive chunking.}
Corrections use 1{,}500-character chunks with 200-character overlap to
preserve exercise boundaries. Course material uses 3{,}000-character chunks
with 250-character overlap to keep theorems intact. The algorithm prefers
natural break points (paragraph breaks, line breaks, sentence endings) over
arbitrary character positions, consistent with findings that structure-aware
chunking improves retrieval~\cite{zhao2024financial_chunking, chen2024mog}.

\paragraph{Metadata extraction.}
Each chunk is annotated with metadata extracted from the GCS path:
\texttt{type} (bac\_officiel, serie, cours), \texttt{chapter},
\texttt{year}, \texttt{session}, \texttt{exo\_id}, \texttt{is\_solution}
(boolean), \texttt{group\_id} (for companion pairing), \texttt{filename},
and \texttt{source} (GCS URI). This metadata enables filtered retrieval and
exercise--correction pairing without manual annotation.

\paragraph{Incremental indexing.}
A manifest file (\texttt{index\_manifest.json}) tracks each file's GCS
generation number and content hash. Only files with changed content are
re-embedded, reducing re-indexing time from minutes to seconds.


% ─────────────────────────────────────────────────────────────────────────────
\subsubsection{Embedding Model Selection}
\label{subsec:embedding-selection}

Retrieval quality is bounded by the embedding model's discriminative power.
If semantically different queries map to similar vectors, no downstream
processing can recover the lost information---especially acute in math,
where minimal notational differences ($v(n)$ vs.\ $u(n)$) carry entirely
different semantics.

\paragraph{Initial experience.}
Our first implementation used a Google embedding model via Vertex AI.
During testing, it did not reliably distinguish between closely related
mathematical sequences---a $v(n)$ query would retrieve $u(n)$ chunks.
General-purpose models treat single-character differences as
insignificant~\cite{mansouri2019tangent, pfahler2020math_embeddings},
but in math they are decisive. This was a qualitative observation, not a
controlled experiment.

\paragraph{Migration to BGE-M3.}
We switched to BGE-M3~\cite{chen2024bge_m3}, built on XLM-RoBERTa with
8{,}192-token context. Its multilingual support (French + Arabic), long
context window, and multi-granularity training make it well-suited for our
heterogeneous corpus. The $v(n)$/$u(n)$ confusion no longer occurred after
migration.

The embedding function wraps the \texttt{BGEM3FlagModel} as a custom
ChromaDB \texttt{EmbeddingFunction}, loaded once as a singleton with
automatic GPU detection.


% ─────────────────────────────────────────────────────────────────────────────
\subsubsection{RAG Retrieval and Generation Engine}
\label{subsec:retrieval-engine}

The runtime pipeline is implemented in \texttt{rag\_engine.py}. It is
UI-agnostic: a single \texttt{.query(question, mode)} method returns a
structured result, callable from Streamlit, Jupyter, or CLI.

\paragraph{Two-stage retrieval.}
The first pass searches correction-type documents
(\texttt{is\_solution=true}), retrieving $k_1 = 10$ candidates. The
rationale is pedagogical: a previously corrected exercise provides not just
mathematical content but a \textit{rhetorical model}---the redaction style
the LLM can mimic.

If the best first-pass distance exceeds 1.2, a second pass retrieves up
to $k_2 = 8$ course-material documents for theorem backing. If both fail,
an unfiltered fallback search is used.

\paragraph{Similarity thresholds.}
Two L2 distance thresholds govern routing:
\begin{description}
    \item[Good match ($\leq 1.2$):] Strong correction match (Case~A). Only
    first-pass results are used.
    \item[Fallback ($\leq 1.6$):] Documents above this are discarded
    entirely.
\end{description}
These were calibrated empirically and are specific to BGE-M3 on our corpus.

\paragraph{Context compilation.}
Selected documents are formatted into XML-tagged context blocks with
metadata, role tags (\texttt{[\'ENONC\'E]} / \texttt{[CORRECTION]}), and
content truncated to 2{,}000 characters per document. Total context is
capped at 14{,}000 characters.

\paragraph{Prompt and generation.}
The system prompt defines the ``Monsieur Tounsi'' persona with:
anti-hallucination rules (use \textit{only} retrieved context), a mimicry
decision tree (Case~A: copy the correction style; Case~B: build from
theorems), a syllabus guard, the Derja sandwich convention, and
anti-injection rules. Generation uses Gemini ($T = 0.15$, max 4{,}096
tokens) with exponential-backoff retry.

\paragraph{Confidence and observability.}
Each query returns timing breakdowns, retrieval metadata, and a confidence
level: \textit{fort} (best distance $\leq 1.2$ and context $>$ 5{,}000
chars), \textit{moyen} (distance $\leq 1.6$), or \textit{faible} (no
useful results).


% ─────────────────────────────────────────────────────────────────────────────
\subsubsection{Paired Retrieval of Exercise and Correction}
\label{subsec:paired-retrieval}

During testing, we found that the system sometimes retrieved a correction
without its exercise statement, or vice versa. A correction is not
self-contained---the notation and constraints in the exercise define what
the correction solves.

\paragraph{Companion fetch mechanism.}
After two-stage retrieval selects its correction documents, the engine
performs a metadata lookup for each: using \texttt{chapter}, \texttt{type},
\texttt{year}, and \texttt{exo\_id}, it queries ChromaDB for chunks with
\texttt{is\_solution=false} (exercise statements). This uses ChromaDB's
direct metadata lookup, not vector search. Up to 3 companion chunks per
correction group are prepended to the context with an
\texttt{[\'ENONC\'E]} label.

This is a specialized form of parent-document
retrieval~\cite{gao2024rag_survey}, adapted to the paired structure of
mathematical exercises. To our knowledge, no prior work has implemented
exercise--correction companion retrieval in a mathematical tutoring system.


% ─────────────────────────────────────────────────────────────────────────────
\subsubsection{Connection to the Tunisian Bac Context}
\label{subsec:bac-context}

Four domain characteristics shaped the system design:

\begin{enumerate}
    \item \textbf{Curriculum boundaries.} The Bac covers a fixed set of ten
    chapters; examiners penalize out-of-program methods. Enforced through
    the syllabus guard and metadata-filtered retrieval.

    \item \textbf{Redaction conventions.} Bac corrections follow a strict
    rhetorical style. The two-stage retrieval prioritizes correction
    exemplars so the LLM can mimic this style.

    \item \textbf{Exercise--correction dependency.} Corrections are always
    read alongside their exercises. The companion fetch preserves this.

    \item \textbf{Bilingual context.} The corpus is in French with \LaTeX{}
    math; student queries may include Tunisian Derja. BGE-M3's multilingual
    training supports this.
\end{enumerate}

% NOTE: If compilation fails, ensure you run pdflatex from the thesis/ directory,
% or adjust the path to ../rag_methods_section if compiling from thesis/chapters/.

% ─────────────────────────────────────────────────────────────────────────────
\subsection{Prompt-Only Baseline}
\label{sec:prompt-only}

The prompt-only baseline serves as a control condition: it uses the same
Gemini model and persona but receives no retrieved documents. Instead, it
encodes the Tunisian Bac mathematics curriculum directly into the system
prompt through ten specialized prompt components:

\begin{enumerate}
    \item \textbf{System persona}: Establishes the ``Monsieur Tounsi''
    identity with curriculum boundaries.

    \item \textbf{Instruction structure}: Defines the expected answer format
    with numbered steps and formal French mathematical phrasing.

    \item \textbf{Reasoning guardrails}: Explicitly forbids methods outside
    the Bac program (e.g., L'H\^{o}pital's rule, Lebesgue integration).

    \item \textbf{Output formatting}: Specifies LaTeX-in-Markdown rendering
    conventions.

    \item \textbf{Hallucination control}: Instructs the model to refuse
    rather than fabricate if unsure about curriculum boundaries.

    \item \textbf{Syllabus guard}: Lists the exact chapters and topics in
    the Bac Math program.

    \item \textbf{Self-verification}: Requires the model to verify each step
    against curriculum constraints before outputting.

    \item \textbf{Self-check block}: Adds an internal consistency check at
    the end of each answer.

    \item \textbf{Refusal policy}: Defines clear refusal messages for
    out-of-scope queries.

    \item \textbf{Confidence calibration}: Maps the model's internal
    certainty to explicit confidence labels.
\end{enumerate}

This baseline tests the hypothesis that a sufficiently detailed prompt can
compensate for the absence of retrieved documents. By comparing it against
the RAG system, we can isolate the value added by retrieval.

% ─────────────────────────────────────────────────────────────────────────────
\subsection{Hybrid Engine}
\label{sec:hybrid}

The hybrid engine combines the strengths of both approaches through a
confidence-based routing mechanism. It operates in three cases:

\begin{itemize}
    \item \textbf{Case~A} (best retrieval distance $\leq 1.2$): The query
    closely matches existing exercises in the corpus. The system uses
    RAG-grounded generation with strong anti-hallucination constraints,
    instructing the model to closely follow the retrieved correction style.

    \item \textbf{Case~B} ($1.2 <$ best distance $\leq 1.6$): The query
    partially matches the corpus. The system combines retrieved context with
    curriculum knowledge from the prompt-only component, allowing the model
    to supplement weak retrievals with parametric knowledge.

    \item \textbf{Case~C} (best distance $> 1.6$ or no results): The query
    has no good match in the corpus. The system falls back entirely to
    prompt-only mode, relying on the encoded curriculum.
\end{itemize}

The distance thresholds (1.2 and 1.6) were determined empirically by
analyzing the distribution of retrieval distances across a sample of typical
student queries. The hybrid engine is implemented as a composition of the
RAG engine (for retrieval) and the prompt-only engine (for curriculum
encoding), without modifying either component.
